\subsection{Noncompletion of the prescribed partial filling}

\begin{proposition}[B.4.30; Proved; the fixed partial filling cannot be completed]
\label{push:B430}
The partial filling of B.4.19 has no proper
Hausdorff complex-analytic completion over its
missing base point. This remains true after finite
coordinate roots of the two base coordinates and
allows arbitrary additional three-dimensional
components of the origin fibre. No K\"ahler or
projective hypothesis is imposed on a putative
completion.

This is a statement about the prescribed partial
filling, including its two punctured axial models.
It does not exclude every possible filling over
the $r$ disc: a different completion along
$r=0$, $q\ne0$ need not have the axial model of
B.4.19. In particular, resolving the graph of the
marked map in B.4.4 does not by itself supply
that additional identification.
\end{proposition}

\begin{proof}
We first prove the local extension statement
needed for traces. We then make a hypothetical
completion equidimensional by an explicit
coefficient-ideal construction and glue multiplied
local traces into the logarithmic torsor
forbidden by B.4.27.

\emph{The analytic logarithmic de Rham input.}
We use \citet[Theorems 6.2 and 6.3,
pp.~32 and 34]{IKN2005E}. For a proper separated
log-smooth analytic map with log-smooth base and
torsion-free cokernels of characteristic groups,
its ordinary relative log-de Rham cohomology is
locally free and, with its Gauss--Manin
connection, is the nilpotent logarithmic
Riemann--Hilbert extension of the log Betti
system. These statements have no K\"ahler or
projectivity assumption. The correction in
\citet{IKN2007ErratumE} concerns the algebraic
Hodge-degeneration Theorem~7.1, which is not
used here. The same distinction is explicit in
\citet[Definition 7.1(i) and Section 9.2,
pp.~164--165 and 173]{KKN2008PicardE}.

\emph{Closed logarithmic forms extend by saturation.}
Consider a proper semistable family of surfaces
\[
 f:Z\longrightarrow\Delta_t\times V,
 \tag{B.4.30.1}
\]
where $V$ is a disc or a point. Its noncentral
fibres are the specified two-dimensional torus
torsors. The semistable charts, with the
transverse parameter retained, are
\[
 t=x_1\cdots x_k,\qquad u=u,\qquad k\le3.
 \tag{B.4.30.2}
\]
Write $\Lambda$ for the integral first-homology
lattice and assume that $T=1+N$, $N^2=0$, and
\[
 \ker(N:\Lambda\longrightarrow\Lambda)
       =\Z w_1\oplus\Z w_2=:W.
 \tag{B.4.30.3}
\]
Here $W$ is the primitive cocharacter lattice
of the specified multiplicative Jacobian
presentation. Suppose its normalized invariant
one-forms $\omega_1,\omega_2$, with periods
$2\pi i\delta_{ij}$ on $w_j$, span a rank-two
subbundle $F_A$ of the canonical extension
$\mathcal H$ of $H^1$.

These assumptions hold at every generic
exceptional divisor in our application. Indeed
$q=t^a$ times a unit and $r=t^b$ times a unit,
with $a,b>0$. The two positive valuation entries
of B.4.27 give (B.4.30.3), with $w_1=w$ and
$w_2=\delta$, after the fixed integral character
change. In the actual period matrix, subtracting
the logarithmic $N$ columns leaves holomorphic
unit periods. The $w,\delta$ columns remain
the identity. Hence their normalized form span
extends as the asserted subbundle $F_A$.

Let $K^\bullet=\Omega^\bullet_{Z/(\Delta\times V)}
(\log)$ and
\[
 \mathcal F=\ker\bigl(f_*\Omega^1(\log)
                   \longrightarrow f_*\Omega^2(\log)\bigr).
\]
This is coherent: the displayed differential
is linear over the base. The triangle
$K^{\ge1}\to K^\bullet\to\OO_Z$ gives the
low-degree exact sequence
\[
 0\longrightarrow\mathcal F
 \longrightarrow\mathcal H
 \longrightarrow R^1f_*\OO_Z.
 \tag{B.4.30.4}
\]
The preceding connecting map is zero since
$f_*\OO_Z=\OO_{\Delta\times V}$ and the
differential is relative. Thus no degeneration
of a Hodge spectral sequence is assumed.

The reduced connected compact central fibres
have only constant holomorphic functions.
For $D=f^{-1}(t=0)$ this gives
$f_*\OO_D=\OO_V$. Applying $f_*$ to
\[
 0\longrightarrow\OO_Z\xrightarrow{\ t\ }
 \OO_Z\longrightarrow\OO_D\longrightarrow0
\]
shows that $R^1f_*\OO_Z$ has no $t$-torsion:
the map $\OO_{\Delta\times V}\to\OO_V$ is
surjective. Consequently $\mathcal F$ is
$t$-saturated in $\mathcal H$. The subbundle
$F_A$ is also $t$-saturated, and the two agree
on $t\ne0$, where all holomorphic one-forms
on a torus are invariant and closed. The
coherent quotients of $\mathcal F+F_A$ by
either summand have support in $t=0$ and
are locally killed by a power of $t$. This
is the localized equality being used; it
does not identify arbitrary punctured
holomorphic sections with meromorphic ones.
Therefore
\[
 \mathcal F=\mathcal H\cap F_A[t^{-1}]=F_A.
 \tag{B.4.30.5}
\]
In particular the normalized forms are actual
closed relative logarithmic forms on $Z$,
holomorphic also in $u$.

\emph{Integral local residues and convergent primitives.}
A local fibre in (B.4.30.2) retracts to
$(S^1)^{k-1}$. Its angular classes are
monodromy-invariant, and their inclusion into
the global fibre intertwines monodromy. By
(B.4.30.3) they belong to $W$. Thus all periods
of $\omega_j/(2\pi i)$ on these angular loops
are constant integers, for all $(t,u)$.

We record the elementary convergent calculation.
Modulo $d\log t$, write a closed form as
\[
 \omega=\sum_{i<k}a_i(x,z,u)d\log x_i
                   +\sum_\ell b_\ell(x,z,u)dz_\ell.
\]
First integrate the ordinary $z$ variables by
the holomorphic homotopy formula. For a
remaining monomial $x^\alpha$, its relative
weight is
\[
 (\alpha_1-\alpha_k,\ldots,\alpha_{k-1}-\alpha_k).
\]
Closedness makes its coefficient vector
proportional to this weight. If the weight is
nonzero, divide by a nonzero integer entry to
integrate the monomial. The denominator has
absolute value at least one, so the series
converges on smaller polydiscs. Zero-weight
monomials are powers of $t$. Their coefficients
are exactly the local angular periods, and
therefore reduce to constant integers in the
normalized forms. We obtain
\[
 \omega_j=d_{\rm rel}h_j+
               \sum_{i<k}n_{ij}d\log x_i,
 \qquad n_{ij}\in\Z,
 \tag{B.4.30.6}
\]
with $h_j$ holomorphic in every variable.
Hence
\[
 \chi_j=e^{h_j}\prod_{i<k}x_i^{n_{ij}}
 \tag{B.4.30.7}
\]
is a meromorphic character whose divisor is
supported on the central coordinate branches.

\emph{Two sections extending into the same cap.}
Let $a,b$ be holomorphic sections of (B.4.30.1).
A finite chain of semistable charts through the
connected compact central fibre joins their
limits. Shrinking $V$ makes the chain uniform
in $u$. Every nonempty open intersection with
a reduced NC fibre contains a smooth point of
that fibre. At such a point $f$ is a
submersion, so choose a holomorphic transverse
auxiliary section in each successive overlap.

In each chart, evaluate the two characters
(B.4.30.7) at the entering and leaving sections
and take their ratios. A pulled-back $x_i$
is a power of $t$ times a holomorphic unit,
since it vanishes only over $t=0$. Thus these
ratios are logarithmic characters on the base.
Multiply them along the finite chain. On each
ordinary fibre, the result is the exponential
of the integrals of $(\omega_1,\omega_2)$ from
$a$ to $b$, so its class modulo the specified
period lattice is exactly $b-a$. A different
path changes it by an integral period. Every
integral ghost is bounded along this full-rank
positive degeneration. Consequently $b-a$
extends as a section of the associated
logarithmic torus $A$.

This construction is holomorphic in the
transverse variable. It is not an inference
from separate one-dimensional restrictions.
It concerns two sections extending into one
cap, and is compatible with B.4.29, where an
essential change of reference origin need not
extend into that cap.

\emph{Equidimensionalization without a projective hypothesis.}
Suppose a proper completion exists. Retain its
reduced main component, normalize it, and make
the finite coordinate roots used in B.4.24.
Its ordinary locus has dimension four; it is
equidimensional of relative dimension two
over the punctured bidisc. Since $r$ is
nonzero on the main component, its origin
fibre has local dimension at most three.

We use the coefficient-ideal argument in
\citet[Corollary 1.2 and its proof,
manuscript pp.~3--5]{Parusinski1990E}, and give
its excess-one specialization. At a point of
the origin fibre with local fibre dimension
three, choose three generic vertical linear
functions $z=(z_1,z_2,z_3)$ so that the germ
\[
 (f,z): (Z,p)\longrightarrow(B\times\C^3,(0,0))
\]
is finite. Its reduced four-dimensional image
is a hypersurface, with an equation
\[
 F(y,z)=\sum_\alpha a_\alpha(y)z^\alpha.
 \tag{B.4.30.8}
\]
The coefficient ideal $I=(a_\alpha)$ is finitely
generated. Moreover $F\in I\OO_{B\times\C^3}$:
reduce modulo $I$ and use injectivity into
the $z$-adic completion, or Krull intersection.
Thus the division below is convergent and
holomorphic, not merely coefficientwise formal.
After shrinking, $I$ is a unit off the origin. Otherwise the image fibre over
some $y\ne0$ would contain an open subset of
$\C^3$; finiteness would force a
three-dimensional source fibre, contrary to
the prescribed partial family.

Blow up $I$. Where its pullback is generated
by $a_{\alpha_0}$, divide (B.4.30.8) by that
generator. The divided equation has a
coefficient equal to one, so is nonzero on
every base fibre. The strict transform lies
in its zero set, and the finite map is
preserved by base change. Its fibres near
this chart therefore have dimension at most
two. A finite number of these charts covers
the compact origin fibre. Blowing up the
product of their coefficient ideals makes
each factor invertible and performs the same
reduction simultaneously. Resolve the
modified surface base. It is dominated by a
finite sequence of point blowups over the
origin. Further base change preserves the
fibre-dimension bound, and a strict transform
cannot increase it. Dominance gives the
opposite bound, so every fibre now has
dimension exactly two.

We only need this equidimensionality, not
flatness of the whole transformed space.
At every point of the modified base it gives
a finite local multisection by a transverse
analytic slice. Normalize that surface. A
normal surface is Cohen--Macaulay, so its
finite map to the smooth surface base is
flat. Thus its positive degree and
multiplicity trace are well defined.

\emph{Uniform semistable reduction along a generic exceptional divisor.}
Let $S\to B$ be the preceding point-blown-up
base and $E$ its compact exceptional tree.
Apply analytic embedded resolution near
the compact exceptional set, in the scope of
\citet[Theorems 1.6 and 1.10]{BM95AB}, to
the total space and the inverse image of $E$
so the latter is \emph{simple} normal crossings. Its components give global
markings of all divisor rays. For each closed
intersection of these components, use its
proper map to a component of $E$. Delete the
critical values of every dominating such map
and the images of nondominating intersections.
These sets are finite. Near any remaining
point, with $u$ a coordinate along $E$, the
map has the form
\[
 t=\text{unit}\prod_{i=1}^k x_i^{m_i},
 \qquad u=u,\qquad k\le3.
 \tag{B.4.30.9}
\]
The last bound follows because the transverse
fibres are surfaces; a fourfold intersection
cannot dominate the curve $E$.

Choose $M$ divisible by the finitely many
multiplicities. Also let $L$ be a common
multiple of $1,\ldots,d_{\max}$ for the
finite multisections to be used. Make the
simultaneous cover $t=\tau^{ML}$ and
normalize. Replace $M$ by $ML$ in the
following lattice calculation. In the lattice
\[
 \sum_i m_i a_i=M b,
\]
each extreme ray has primitive generator
$a=(M/m_i)e_i$, $b=1$. Thus the height-one
slice is an integral polygon of dimension
at most two. Insert all its lattice points
by star subdivisions, using one order on the
finite marked cone complex. Shared faces
receive the same subdivisions. The final
empty lattice triangles are unimodular, so
the toric modification is smooth and its
$\tau$ fibre is reduced NC. The monomial
ideals glue under coordinate changes by
units. This is precisely the analytic
subdivision construction of
\citet[Lemma 2.14, pp.~6--7]{EH2024E}; their
projective main theorem is not used.
The construction retains the transverse
parameter $u$ throughout.

The finite multisections split on this
simultaneous cyclic cover at generic points
of $E$.
Indeed, away from the remaining branch-curve
intersections, their finite normal covers
are the cyclic covers of a punctured disc,
times the $u$ disc. A proper modification
gives a meromorphic lift of each resulting
section; on its normal surface base the
indeterminacy is isolated. Delete those
finitely many additional values of $u$.
The sections then lift holomorphically into
the semistable model. The preceding
saturation and primitive argument applies.

If a difference extends only after
$t=\tau^e$, take its multiplicative norm.
For each character the formula is
\[
 \prod_{j=0}^{e-1}(\zeta_e^j\tau)^k
                      U(\zeta_e^j\tau,u)
   =(-1)^{k(e-1)}t^k\,\text{unit}.
 \tag{B.4.30.10}
\]
A change of lift adds an integral period,
whose norm remains a period. The resulting
section downstairs is $e$ times the original
difference. This is an explicit norm on the
usual analytic site, not an assumption of
Kummer descent for a quotient sheaf. A
common multiplier suffices: there are
finitely many divisor multiplicities and
multisections, and ramification indices of
a finite multisection are bounded by its
degree. The product of the multiplicity
common multiple and
$\operatorname{lcm}(1,\ldots,d_{\max})$
bounds the simultaneous root and norm.
The toric subdivision is performed after
this combined root, so no further base
change silently destroys semistability.

\emph{Hartogs for logarithmic sections, including the period ambiguity.}
Let $V_p$ be a small bidisc about an isolated
remaining point on $S$, and put
$V_p^*=V_p\setminus\{p\}$. The period sheaf
$Y=\Z^2$ is constant and
\[
 H^1(V_p^*,Y)=H^1(S^3,\Z^2)=0.
\]
The usual-site exact sequence
\[
 0\longrightarrow Y\longrightarrow
 T_{\log}^{(Y)}\longrightarrow A\longrightarrow0
 \tag{B.4.30.11}
\]
therefore lifts every $A$-section on $V_p^*$
to bounded logarithmic characters. By
B.4.27 the middle sheaf is exactly
\[
 M_{D_Q+E}^{\rm gp}\times
 M_{D_Q+D_R+E}^{\rm gp}.
\]
For either character divide by the local
boundary monomial carrying its integer
orders. The resulting function and its
inverse are holomorphic on the punctured
bidisc. Hartogs extends both, and their
product remains one. Restore the monomial.
The forbidden strict-$r$ order of the first
character stays zero, so the extension is
still bounded.

The extension is unique. Locally lift the
difference of two extensions to the middle
sheaf. On the connected punctured bidisc
it is an integral period, hence one constant
element of $Y$. Subtract that period and
apply the identity theorem to its two
characters. This proves uniqueness and
preserves all cocycle identities. It does
not assert Hartogs extension for $A$-torsors.

\emph{Multiplying local traces and gluing the forbidden torsor.}
Choose a finite cover $U_i$ of a
neighbourhood of the compact exceptional
tree and a finite local multisection of
degree $d_i>0$ on each member. Let $d$ be
a common positive multiple of these degrees.
On $U_i\setminus E$, its trace multiplied
by $d/d_i$ is a section $t_i$ of the
$d$-fold contracted product $\mathcal P^d$
of the actual punctured $A$-torsor.
On the two original punctured axes these
are genuine logarithmic trace sections:
a holomorphic map into either prescribed
fan model pulls its characters back to
powers of the local root coordinate times
units, and their norms are exactly those
in B.4.26. A proper modification projecting
to that model does not change this argument.

On overlaps, $t_i-t_j$ is an $A$-section
away from $E$. The uniform argument above
gives a positive multiple extending across
the generic part of every exceptional
component. Choose one common multiplier
$m$ for the finite set of curves and
multisections and replace $t_i$ by $mt_i$.
These are sections of $\mathcal P^{dm}$.
Their differences now extend away from
finitely many points of $E$; the Hartogs
argument extends them across those points
as well, including crossings. Uniqueness
preserves the cocycle equations.

Glue trivial $A$-torsors on the $U_i$ by
these extended differences and glue them
to $\mathcal P^{dm}$ off $E$ using the
trace identifications. This produces an
$A$-torsor on $S$ restricting to
$\mathcal P^{dm}$ on
$\beta^{-1}(B\setminus\{0\})$.
B.4.27 says that $H^1(S,A)=0$.
But B.4.24 says that every positive
multiple of this actual punctured torsor
is nontrivial: its $r$-loop connecting
class is a nonzero multiple of $\gamma$.
This is the required contradiction.

Finally, all period and cohomology arguments
in B.4.24 and B.4.27 allow further finite
coordinate roots. The proof is therefore
unchanged by them. Its use of the two
\emph{prescribed} punctured axial models
is essential. Along the strict $r$ axis
alone the monodromy has rank one, so an
arbitrary other cap can have additional
invariant elliptic residues. The proof
does not silently identify such a cap
with B.4.19.
\end{proof}
